请注意，这里的地址是相对于系统寄存器基地址的地址偏移量（相对地址）。请参阅章节 \ref{subsection:system-reg-base-address} 获取有关系统寄存器基地址的信息。

\begin{register}{H}{SYSTEM\_ROM\_CTRL\_0\_REG}{\hexprefix{}0000}\label{regdesc:SYSTEMROMCTRL0REG}
\regfield{(reserved)}{30}{2}{000000000000000000000000000000}%
\regfield{SYSTEM\_ROM\_FO}{2}{0}{{\hexprefix{}3}}%
\reglabel{Reset}\regnewline%
\begin{regdesc}\begin{reglist}
\label{fielddesc:SYSTEMROMFO}\item [SYSTEM\_ROM\_FO] 强制打开内部 ROM 的时钟门控。详见表 \ref{table:romcontrollingbit}。（读/写）

\end{reglist}\end{regdesc}
\end{register}

\begin{register}{H}{SYSTEM\_ROM\_CTRL\_1\_REG}{\hexprefix{}0004}\label{regdesc:SYSTEMROMCTRL1REG}
\regfield{(reserved)}{28}{4}{0000000000000000000000000000}%
\regfield{SYSTEM\_ROM\_FORCE\_PU}{2}{2}{{3}}%
\regfield{SYSTEM\_ROM\_FORCE\_PD}{2}{0}{{0}}%
\reglabel{Reset}\regnewline%
\begin{regdesc}\begin{reglist}
\label{fielddesc:SYSTEMROMFORCEPD}\item [SYSTEM\_ROM\_FORCE\_PD] 控制内部 ROM 的掉电。详见表 \ref{table:romcontrollingbit}。（读/写）
\label{fielddesc:SYSTEMROMFORCEPU}\item [SYSTEM\_ROM\_FORCE\_PU] 控制内部 ROM 的上电。详见表 \ref{table:romcontrollingbit}。（读/写）

\end{reglist}\end{regdesc}
\end{register}


\begin{register}{H}{SYSTEM\_SRAM\_CTRL\_0\_REG}{\hexprefix{}0008}\label{regdesc:SYSTEMSRAMCTRL0REG}
\regfield{(reserved)}{10}{22}{0000000000}%
\regfield{SYSTEM\_SRAM\_FO}{22}{0}{{\hexprefix{}3fffff}}%
\reglabel{Reset}\regnewline%
\begin{regdesc}\begin{reglist}
\label{fielddesc:SYSTEMSRAMFO}\item [SYSTEM\_SRAM\_FO] 强制打开内部 SRAM 的时钟门控。详见表 \ref{table:sramcontrollingbit}。（读/写）

\end{reglist}\end{regdesc}
\end{register}

\begin{register}{H}{SYSTEM\_SRAM\_CTRL\_1\_REG}{\hexprefix{}000C}\label{regdesc:SYSTEMSRAMCTRL1REG}
\regfield{(reserved)}{10}{22}{0000000000}%
\regfield{SYSTEM\_SRAM\_FORCE\_PD}{22}{0}{{0}}%
\reglabel{Reset}\regnewline%
\begin{regdesc}\begin{reglist}
\label{fielddesc:SYSTEMSRAMFORCEPD}\item [SYSTEM\_SRAM\_FORCE\_PD] 控制内部 SRAM 的掉电。详见表 \ref{table:sramcontrollingbit}。（读/写）

\end{reglist}\end{regdesc}
\end{register}



\begin{register}{H}{SYSTEM\_CPU\_PERI\_CLK\_EN\_REG}{\hexprefix{}0010}\label{regdesc:SYSTEMCPUPERICLKENREG}
\regfield{(reserved)}{24}{8}{000000000000000000000000}%
\regfield{SYSTEM\_CLK\_EN\_DEDICATED\_GPIO}{1}{7}{{0}}%
\regfield{(reserved)}{1}{6}{{0}}%
\regfield{(reserved)}{6}{0}{000000}%
\reglabel{Reset}\regnewline%
\begin{regdesc}\begin{reglist}
\label{fielddesc:SYSTEMCLKENDEDICATEDGPIO}\item [SYSTEM\_CLK\_EN\_DEDICATED\_GPIO] 置 1 打开 DEDICATED GPIO 时钟。详见章节 \ref{mod:iomuxgpio} \textit{\nameref{mod:iomuxgpio}}。（读/写）
\end{reglist}\end{regdesc}
\end{register}


\begin{register}{H}{SYSTEM\_CPU\_PERI\_RST\_EN\_REG}{\hexprefix{}0014}\label{regdesc:SYSTEMCPUPERIRSTENREG}
\regfield{(reserved)}{24}{8}{000000000000000000000000}%
\regfield{SYSTEM\_RST\_EN\_DEDICATED\_GPIO}{1}{7}{{1}}%
\regfield{(reserved)}{1}{6}{{1}}%
\regfield{(reserved)}{6}{0}{000000}%
\reglabel{Reset}\regnewline%
\begin{regdesc}\begin{reglist}
\label{fielddesc:SYSTEMRSTENDEDICATEDGPIO}\item [SYSTEM\_RST\_EN\_DEDICATED\_GPIO] 置 1 复位 DEDICATED GPIO。详见章节 \ref{mod:iomuxgpio} \textit{\nameref{mod:iomuxgpio}}。（读/写）
\end{reglist}\end{regdesc}
\end{register}


\begin{register}{H}{SYSTEM\_CPU\_PER\_CONF\_REG}{\hexprefix{}0018}\label{regdesc:SYSTEMCPUPERCONFREG}
\regfield{(reserved)}{24}{8}{000000000000000000000000}%
\regfield{SYSTEM\_CPU\_WAITI\_DELAY\_NUM}{4}{4}{{\hexprefix{}0}}%
\regfield{SYSTEM\_CPU\_WAIT\_MODE\_FORCE\_ON}{1}{3}{{1}}%
\regfield{SYSTEM\_PLL\_FREQ\_SEL}{1}{2}{{1}}%
\regfield{SYSTEM\_CPUPERIOD\_SEL}{2}{0}{{0}}%
\reglabel{Reset}\regnewline%
\begin{regdesc}\begin{reglist}
\label{fielddesc:SYSTEMCPUPERIODSEL}\item [SYSTEM\_CPUPERIOD\_SEL] 选择 CPU 时钟周期或时钟频率。详见章节 \ref{mod:resclk} \textit{\nameref{mod:resclk}}中表 \ref{table:cpu clksrc}。（读/写）
\label{fielddesc:SYSTEMPLLFREQSEL}\item [SYSTEM\_PLL\_FREQ\_SEL] 选择基于 CPU 时钟周期的 PPL 时钟频率。详见章节 \ref{mod:resclk} \textit{\nameref{mod:resclk}}。（读/写）
\label{fielddesc:SYSTEMCPUWAITMODEFORCEON}\item [SYSTEM\_CPU\_WAIT\_MODE\_FORCE\_ON] 置 1 强制打开 CPU 等待模式。在 CPU 等待模式下，CPU 时钟门控一直处于关闭状态，直到中断产生。用户也可通过 WAITI 指令强制打开 CPU 等待模式。（读/写）

\label{fielddesc:SYSTEMCPUWAITIDELAYNUM}\item [SYSTEM\_CPU\_WAITI\_DELAY\_NUM] 设置 CPU 在收到 WAITI 指令后进入 CPU 等待模式前的等待周期。（读/写）

\end{reglist}\end{regdesc}
\end{register}


\begin{register}{H}{SYSTEM\_JTAG\_CTRL\_0\_REG}{\hexprefix{}001C}\label{regdesc:SYSTEMJTAGCTRL0REG}
\regfield{SYSTEM\_CANCEL\_EFUSE\_DISABLE\_JTAG\_TEMPORARY\_0}{32}{0}{{0}}%
\reglabel{Reset}\regnewline%
\begin{regdesc}\begin{reglist}
\label{fielddesc:SYSTEMCANCELEFUSEDISABLEJTAGTEMPORARY0}\item [SYSTEM\_CANCEL\_EFUSE\_DISABLE\_JTAG\_TEMPORARY\_0] 本组寄存器（共 256 位）用于撤销 eFuse 对 JTAG 的临时关闭，本寄存器为第 0 到 31 位。具体请见章节 \ref{mod:hmac} \textit{\nameref{mod:hmac}}。（只写）
\end{reglist}\end{regdesc}
\end{register}


\begin{register}{H}{SYSTEM\_JTAG\_CTRL\_1\_REG}{\hexprefix{}0020}\label{regdesc:SYSTEMJTAGCTRL1REG}
\regfield{SYSTEM\_CANCEL\_EFUSE\_DISABLE\_JTAG\_TEMPORARY\_1}{32}{0}{{0}}%
\reglabel{Reset}\regnewline%
\begin{regdesc}\begin{reglist}
\label{fielddesc:SYSTEMCANCELEFUSEDISABLEJTAGTEMPORARY1}\item [SYSTEM\_CANCEL\_EFUSE\_DISABLE\_JTAG\_TEMPORARY\_1] 本组寄存器（共 256 位）用于撤销 eFuse 对 JTAG 的临时关闭，本寄存器为第 32 到 63 位。具体请见章节 \ref{mod:hmac} \textit{\nameref{mod:hmac}}。（只写）
\end{reglist}\end{regdesc}
\end{register}


\begin{register}{H}{SYSTEM\_JTAG\_CTRL\_2\_REG}{\hexprefix{}0024}\label{regdesc:SYSTEMJTAGCTRL2REG}
\regfield{SYSTEM\_CANCEL\_EFUSE\_DISABLE\_JTAG\_TEMPORARY\_2}{32}{0}{{0}}%
\reglabel{Reset}\regnewline%
\begin{regdesc}\begin{reglist}
\label{fielddesc:SYSTEMCANCELEFUSEDISABLEJTAGTEMPORARY2}\item [SYSTEM\_CANCEL\_EFUSE\_DISABLE\_JTAG\_TEMPORARY\_2] 本组寄存器（共 256 位）用于撤销 eFuse 对 JTAG 的临时关闭，本寄存器为第 64 到 95 位。具体请见章节 \ref{mod:hmac} \textit{\nameref{mod:hmac}}。（只写）
\end{reglist}\end{regdesc}
\end{register}


\begin{register}{H}{SYSTEM\_JTAG\_CTRL\_3\_REG}{\hexprefix{}0028}\label{regdesc:SYSTEMJTAGCTRL3REG}
\regfield{SYSTEM\_CANCEL\_EFUSE\_DISABLE\_JTAG\_TEMPORARY\_3}{32}{0}{{0}}%
\reglabel{Reset}\regnewline%
\begin{regdesc}\begin{reglist}
\label{fielddesc:SYSTEMCANCELEFUSEDISABLEJTAGTEMPORARY3}\item [SYSTEM\_CANCEL\_EFUSE\_DISABLE\_JTAG\_TEMPORARY\_3] 本组寄存器（共 256 位）用于撤销 eFuse 对 JTAG 的临时关闭，本寄存器为第 96 到 127 位。具体请见章节 \ref{mod:hmac} \textit{\nameref{mod:hmac}}。（只写）
\end{reglist}\end{regdesc}
\end{register}


\begin{register}{H}{SYSTEM\_JTAG\_CTRL\_4\_REG}{\hexprefix{}002C}\label{regdesc:SYSTEMJTAGCTRL4REG}
\regfield{SYSTEM\_CANCEL\_EFUSE\_DISABLE\_JTAG\_TEMPORARY\_4}{32}{0}{{0}}%
\reglabel{Reset}\regnewline%
\begin{regdesc}\begin{reglist}
\label{fielddesc:SYSTEMCANCELEFUSEDISABLEJTAGTEMPORARY4}\item [SYSTEM\_CANCEL\_EFUSE\_DISABLE\_JTAG\_TEMPORARY\_4] 本组寄存器（共 256 位）用于撤销 eFuse 对 JTAG 的临时关闭，本寄存器为第 128 到 159 位。具体请见章节 \ref{mod:hmac} \textit{\nameref{mod:hmac}}。（只写）
\end{reglist}\end{regdesc}
\end{register}


\begin{register}{H}{SYSTEM\_JTAG\_CTRL\_5\_REG}{\hexprefix{}0030}\label{regdesc:SYSTEMJTAGCTRL5REG}
\regfield{SYSTEM\_CANCEL\_EFUSE\_DISABLE\_JTAG\_TEMPORARY\_5}{32}{0}{{0}}%
\reglabel{Reset}\regnewline%
\begin{regdesc}\begin{reglist}
\label{fielddesc:SYSTEMCANCELEFUSEDISABLEJTAGTEMPORARY5}\item [SYSTEM\_CANCEL\_EFUSE\_DISABLE\_JTAG\_TEMPORARY\_5] 本组寄存器（共 256 位）用于撤销 eFuse 对 JTAG 的临时关闭，本寄存器为第 160 到 191 位。具体请见章节 \ref{mod:hmac} \textit{\nameref{mod:hmac}}。（只写）
\end{reglist}\end{regdesc}
\end{register}


\begin{register}{H}{SYSTEM\_JTAG\_CTRL\_6\_REG}{\hexprefix{}0034}\label{regdesc:SYSTEMJTAGCTRL6REG}
\regfield{SYSTEM\_CANCEL\_EFUSE\_DISABLE\_JTAG\_TEMPORARY\_6}{32}{0}{{0}}%
\reglabel{Reset}\regnewline%
\begin{regdesc}\begin{reglist}
\label{fielddesc:SYSTEMCANCELEFUSEDISABLEJTAGTEMPORARY6}\item [SYSTEM\_CANCEL\_EFUSE\_DISABLE\_JTAG\_TEMPORARY\_6] 本组寄存器（共 256 位）用于撤销 eFuse 对 JTAG 的临时关闭，本寄存器为第 192 到 223 位。具体请见第 \ref{mod:hmac} 章：\textit{\nameref{mod:hmac}}。（只写）
\end{reglist}\end{regdesc}
\end{register}


\begin{register}{H}{SYSTEM\_JTAG\_CTRL\_7\_REG}{\hexprefix{}0038}\label{regdesc:SYSTEMJTAGCTRL7REG}
\regfield{SYSTEM\_CANCEL\_EFUSE\_DISABLE\_JTAG\_TEMPORARY\_7}{32}{0}{{0}}%
\reglabel{Reset}\regnewline%
\begin{regdesc}\begin{reglist}
\label{fielddesc:SYSTEMCANCELEFUSEDISABLEJTAGTEMPORARY7}\item [SYSTEM\_CANCEL\_EFUSE\_DISABLE\_JTAG\_TEMPORARY\_7] 本组寄存器（共 256 位）用于撤销 eFuse 对 JTAG 的临时关闭，本寄存器为第 224 到 255 位。具体请见章节 \ref{mod:hmac} \textit{\nameref{mod:hmac}}。（只写）
\end{reglist}\end{regdesc}
\end{register}


\begin{register}{H}{SYSTEM\_MEM\_PD\_MASK\_REG}{\hexprefix{}003C}\label{regdesc:SYSTEMMEMPDMASKREG}
\regfield{(reserved)}{31}{1}{0000000000000000000000000000000}%
\regfield{SYSTEM\_LSLP\_MEM\_PD\_MASK}{1}{0}{{1}}%
\reglabel{Reset}\regnewline%
\begin{regdesc}\begin{reglist}
\label{fielddesc:SYSTEMLSLPMEMPDMASK}\item [SYSTEM\_LSLP\_MEM\_PD\_MASK] 置 1 允许存储器在 light-sleep 模式下仍正常工作。（读/写）
\end{reglist}\end{regdesc}
\end{register}


\begin{register}{H}{SYSTEM\_PERIP\_CLK\_EN0\_REG}{\hexprefix{}0040}\label{regdesc:SYSTEMPERIPCLKEN0REG}
\regfield{(reserved)}{1}{31}{{1}}%
\regfield{SYSTEM\_ADC2\_ARB\_CLK\_EN}{1}{30}{{1}}%
\regfield{SYSTEM\_SYSTIMER\_CLK\_EN}{1}{29}{{1}}%
\regfield{SYSTEM\_APB\_SARADC\_CLK\_EN}{1}{28}{{1}}%
\regfield{SYSTEM\_SPI3\_DMA\_CLK\_EN}{1}{27}{{1}}%
\regfield{SYSTEM\_PWM3\_CLK\_EN}{1}{26}{{0}}%
\regfield{SYSTEM\_PWM2\_CLK\_EN}{1}{25}{{0}}%
\regfield{SYSTEM\_UART\_MEM\_CLK\_EN}{1}{24}{{1}}%
\regfield{SYSTEM\_USB\_CLK\_EN}{1}{23}{{1}}%
\regfield{SYSTEM\_SPI2\_DMA\_CLK\_EN}{1}{22}{{1}}%
\regfield{(reserved)}{1}{21}{{0}}%
\regfield{SYSTEM\_PWM1\_CLK\_EN}{1}{20}{{0}}%
\regfield{SYSTEM\_CAN\_CLK\_EN}{1}{19}{{0}}%
\regfield{SYSTEM\_I2C\_EXT1\_CLK\_EN}{1}{18}{{0}}%
\regfield{SYSTEM\_PWM0\_CLK\_EN}{1}{17}{{0}}%
\regfield{SYSTEM\_SPI3\_CLK\_EN}{1}{16}{{1}}%
\regfield{SYSTEM\_TIMERGROUP1\_CLK\_EN}{1}{15}{{1}}%
\regfield{SYSTEM\_EFUSE\_CLK\_EN}{1}{14}{{1}}%
\regfield{SYSTEM\_TIMERGROUP\_CLK\_EN}{1}{13}{{1}}%
\regfield{SYSTEM\_UHCI1\_CLK\_EN}{1}{12}{{0}}%
\regfield{SYSTEM\_LEDC\_CLK\_EN}{1}{11}{{0}}%
\regfield{SYSTEM\_PCNT\_CLK\_EN}{1}{10}{{0}}%
\regfield{SYSTEM\_RMT\_CLK\_EN}{1}{9}{{0}}%
\regfield{SYSTEM\_UHCI0\_CLK\_EN}{1}{8}{{0}}%
\regfield{SYSTEM\_I2C\_EXT0\_CLK\_EN}{1}{7}{{0}}%
\regfield{SYSTEM\_SPI2\_CLK\_EN}{1}{6}{{1}}%
\regfield{SYSTEM\_UART1\_CLK\_EN}{1}{5}{{1}}%
\regfield{SYSTEM\_I2S0\_CLK\_EN}{1}{4}{{0}}%
\regfield{SYSTEM\_WDG\_CLK\_EN}{1}{3}{{1}}%
\regfield{SYSTEM\_UART\_CLK\_EN}{1}{2}{{1}}%
\regfield{SYSTEM\_SPI01\_CLK\_EN}{1}{1}{{1}}%
\regfield{SYSTEM\_TIMERS\_CLK\_EN}{1}{0}{{1}}%
\reglabel{Reset}\regnewline%
\begin{regdesc}\begin{reglist}
\item[SYSTEM\_CPU\_PERI\_CLK\_EN0\_REG] 使能不同的外设时钟，详见表 \ref{table:peripheralcontrollingbit}。

\end{reglist}\end{regdesc}
\end{register}




\begin{register}{H}{SYSTEM\_PERIP\_CLK\_EN1\_REG}{\hexprefix{}0044}\label{regdesc:SYSTEMPERIPCLKEN1REG}
\regfield{(reserved)}{25}{7}{0000000000000000000000000}%
\regfield{SYSTEM\_CRYPTO\_DMA\_CLK\_EN}{1}{6}{{0}}%
\regfield{SYSTEM\_CRYPTO\_HMAC\_CLK\_EN}{1}{5}{{0}}%
\regfield{SYSTEM\_CRYPTO\_DS\_CLK\_EN}{1}{4}{{0}}%
\regfield{SYSTEM\_CRYPTO\_RSA\_CLK\_EN}{1}{3}{{0}}%
\regfield{SYSTEM\_CRYPTO\_SHA\_CLK\_EN}{1}{2}{{0}}%
\regfield{SYSTEM\_CRYPTO\_AES\_CLK\_EN}{1}{1}{{0}}%
\regfield{(reserved)}{1}{0}{0}%
\reglabel{Reset}\regnewline%
\begin{regdesc}\begin{reglist}
\label{fielddesc:SYSTEMCRYPTOAESCLKEN}\item [SYSTEM\_CRYPTO\_AES\_CLK\_EN] 置 1 打开 AES 加速器的时钟。（读/写）
\label{fielddesc:SYSTEMCRYPTOSHACLKEN}\item [SYSTEM\_CRYPTO\_SHA\_CLK\_EN] 置 1 打开 SHA 加速器的时钟。（读/写）
\label{fielddesc:SYSTEMCRYPTORSACLKEN}\item [SYSTEM\_CRYPTO\_RSA\_CLK\_EN] 置 1 打开 RSA 加速器的时钟。（读/写）
\label{fielddesc:SYSTEMCRYPTODSCLKEN}\item [SYSTEM\_CRYPTO\_DS\_CLK\_EN] 置 1 打开 RSA\_DS 外设的时钟。（读/写）
\label{fielddesc:SYSTEMCRYPTOHMACCLKEN}\item [SYSTEM\_CRYPTO\_HMAC\_CLK\_EN] 置 1 打开 HMAC 加速器的时钟。（读/写）
\label{fielddesc:SYSTEMCRYPTODMACLKEN}\item [SYSTEM\_CRYPTO\_DMA\_CLK\_EN] 置 1 打开加密 DMA 的时钟。（读/写）
\end{reglist}\end{regdesc}
\end{register}


\begin{register}{H}{SYSTEM\_PERIP\_RST\_EN0\_REG}{\hexprefix{}0048}\label{regdesc:SYSTEMPERIPRSTEN0REG}
\regfield{(reserved)}{1}{31}{{0}}%
\regfield{SYSTEM\_ADC2\_ARB\_RST}{1}{30}{{0}}%
\regfield{SYSTEM\_SYSTIMER\_RST}{1}{29}{{0}}%
\regfield{SYSTEM\_APB\_SARADC\_RST}{1}{28}{{0}}%
\regfield{SYSTEM\_SPI3\_DMA\_RST}{1}{27}{{0}}%
\regfield{SYSTEM\_PWM3\_RST}{1}{26}{{0}}%
\regfield{SYSTEM\_PWM2\_RST}{1}{25}{{0}}%
\regfield{SYSTEM\_UART\_MEM\_RST}{1}{24}{{0}}%
\regfield{SYSTEM\_USB\_RST}{1}{23}{{0}}%
\regfield{SYSTEM\_SPI2\_DMA\_RST}{1}{22}{{0}}%
\regfield{(reserved)}{1}{21}{{0}}%
\regfield{SYSTEM\_PWM1\_RST}{1}{20}{{0}}%
\regfield{SYSTEM\_CAN\_RST}{1}{19}{{0}}%
\regfield{SYSTEM\_I2C\_EXT1\_RST}{1}{18}{{0}}%
\regfield{SYSTEM\_PWM0\_RST}{1}{17}{{0}}%
\regfield{SYSTEM\_SPI3\_RST}{1}{16}{{0}}%
\regfield{SYSTEM\_TIMERGROUP1\_RST}{1}{15}{{0}}%
\regfield{SYSTEM\_EFUSE\_RST}{1}{14}{{0}}%
\regfield{SYSTEM\_TIMERGROUP\_RST}{1}{13}{{0}}%
\regfield{SYSTEM\_UHCI1\_RST}{1}{12}{{0}}%
\regfield{SYSTEM\_LEDC\_RST}{1}{11}{{0}}%
\regfield{SYSTEM\_PCNT\_RST}{1}{10}{{0}}%
\regfield{SYSTEM\_RMT\_RST}{1}{9}{{0}}%
\regfield{SYSTEM\_UHCI0\_RST}{1}{8}{{0}}%
\regfield{SYSTEM\_I2C\_EXT0\_RST}{1}{7}{{0}}%
\regfield{SYSTEM\_SPI2\_RST}{1}{6}{{0}}%
\regfield{SYSTEM\_UART1\_RST}{1}{5}{{0}}%
\regfield{SYSTEM\_I2S0\_RST}{1}{4}{{0}}%
\regfield{SYSTEM\_WDG\_RST}{1}{3}{{0}}%
\regfield{SYSTEM\_UART\_RST}{1}{2}{{0}}%
\regfield{SYSTEM\_SPI01\_RST}{1}{1}{{0}}%
\regfield{SYSTEM\_TIMERS\_RST}{1}{0}{{0}}%
\reglabel{Reset}\regnewline%
\begin{regdesc}\begin{reglist}
\item[SYSTEM\_PERIP\_RST\_EN0\_REG] 复位不同外设，详见表 \ref{table:peripheralcontrollingbit}。


\end{reglist}\end{regdesc}
\end{register}


\begin{register}{H}{SYSTEM\_PERIP\_RST\_EN1\_REG}{\hexprefix{}004C}\label{regdesc:SYSTEMPERIPRSTEN1REG}
\regfield{(reserved)}{25}{7}{0000000000000000000000000}%
\regfield{SYSTEM\_CRYPTO\_DMA\_RST}{1}{6}{{1}}%
\regfield{SYSTEM\_CRYPTO\_HMAC\_RST}{1}{5}{{1}}%
\regfield{SYSTEM\_CRYPTO\_DS\_RST}{1}{4}{{1}}%
\regfield{SYSTEM\_CRYPTO\_RSA\_RST}{1}{3}{{1}}%
\regfield{SYSTEM\_CRYPTO\_SHA\_RST}{1}{2}{{1}}%
\regfield{SYSTEM\_CRYPTO\_AES\_RST}{1}{1}{{1}}%
\regfield{(reserved)}{1}{0}{0}%
\reglabel{Reset}\regnewline%
\begin{regdesc}\begin{reglist}
\item[SYSTEM\_PERIP\_RST\_EN1\_REG] 复位不同加速器，详见表 \ref{table:peripheralcontrollingbit}。

\end{reglist}\end{regdesc}
\end{register}








\begin{register}{H}{SYSTEM\_BT\_LPCK\_DIV\_FRAC\_REG}{\hexprefix{}0054}\label{regdesc:SYSTEMBTLPCKDIVFRACREG}
\regfield{(reserved)}{3}{29}{000}%
\regfield{SYSTEM\_LPCLK\_RTC\_EN}{1}{28}{{0}}%
\regfield{SYSTEM\_LPCLK\_SEL\_XTAL32K}{1}{27}{{0}}%
\regfield{SYSTEM\_LPCLK\_SEL\_XTAL}{1}{26}{{0}}%
\regfield{SYSTEM\_LPCLK\_SEL\_8M}{1}{25}{{1}}%
\regfield{SYSTEM\_LPCLK\_SEL\_RTC\_SLOW}{1}{24}{{0}}%
\regfield{(reserved)}{12}{12}{{1}}%
\regfield{(reserved)}{12}{0}{{1}}%
\reglabel{Reset}\regnewline%
\begin{regdesc}\begin{reglist}
\label{fielddesc:SYSTEMLPCLKSELRTCSLOW}\item [SYSTEM\_LPCLK\_SEL\_RTC\_SLOW] 置 1 选择 RTC\_SLOW\_CLK 为低功耗时钟。（读/写）
\label{fielddesc:SYSTEMLPCLKSEL8M}\item [SYSTEM\_LPCLK\_SEL\_8M] 置 1 选择 RC\_FAST\_CLK 为低功耗时钟。（读/写）
\label{fielddesc:SYSTEMLPCLKSELXTAL}\item [SYSTEM\_LPCLK\_SEL\_XTAL] 置 1 选择 XTAL\_CLK 为低功耗时钟。（读/写）
\label{fielddesc:SYSTEMLPCLKSELXTAL32K}\item [SYSTEM\_LPCLK\_SEL\_XTAL32K] 置 1 选择 XTAL32K\_CLK 为低功耗时钟。（读/写）
\label{fielddesc:SYSTEMLPCLKRTCEN}\item [SYSTEM\_LPCLK\_RTC\_EN] 置 1 使能 RTC 低功耗时钟。（读/写）
\end{reglist}\end{regdesc}
\end{register}



\begin{register}{H}{SYSTEM\_CPU\_INTR\_FROM\_CPU\_0\_REG}{\hexprefix{}0058}\label{regdesc:SYSTEMCPUINTRFROMCPU0REG}
\regfield{(reserved)}{31}{1}{0000000000000000000000000000000}%
\regfield{SYSTEM\_CPU\_INTR\_FROM\_CPU\_0}{1}{0}{{0}}%
\reglabel{Reset}\regnewline%
\begin{regdesc}\begin{reglist}
\label{fielddesc:SYSTEMCPUINTRFROMCPU0}\item [SYSTEM\_CPU\_INTR\_FROM\_CPU\_0]置 1 生成 CPU 中断 0。该位需在 ISR 过程中由软件清 0。（读/写）
\end{reglist}\end{regdesc}
\end{register}

\begin{register}{H}{SYSTEM\_CPU\_INTR\_FROM\_CPU\_1\_REG}{\hexprefix{}005C}\label{regdesc:SYSTEMCPUINTRFROMCPU1REG}
\regfield{(reserved)}{31}{1}{0000000000000000000000000000000}%
\regfield{SYSTEM\_CPU\_INTR\_FROM\_CPU\_1}{1}{0}{{0}}%
\reglabel{Reset}\regnewline%
\begin{regdesc}\begin{reglist}
\label{fielddesc:SYSTEMCPUINTRFROMCPU1}\item [SYSTEM\_CPU\_INTR\_FROM\_CPU\_1] 置 1 生成 CPU 中断 1。该位需在 ISR 过程中由软件清 0。（读/写）
\end{reglist}\end{regdesc}
\end{register}

\begin{register}{H}{SYSTEM\_CPU\_INTR\_FROM\_CPU\_2\_REG}{\hexprefix{}0060}\label{regdesc:SYSTEMCPUINTRFROMCPU2REG}
\regfield{(reserved)}{31}{1}{0000000000000000000000000000000}%
\regfield{SYSTEM\_CPU\_INTR\_FROM\_CPU\_2}{1}{0}{{0}}%
\reglabel{Reset}\regnewline%
\begin{regdesc}\begin{reglist}
\label{fielddesc:SYSTEMCPUINTRFROMCPU2}\item [SYSTEM\_CPU\_INTR\_FROM\_CPU\_2] 置 1 生成 CPU 中断 2。该位需在 ISR 过程中由软件清 0。（读/写）
\end{reglist}\end{regdesc}
\end{register}

\begin{register}{H}{SYSTEM\_CPU\_INTR\_FROM\_CPU\_3\_REG}{\hexprefix{}0064}\label{regdesc:SYSTEMCPUINTRFROMCPU3REG}
\regfield{(reserved)}{31}{1}{0000000000000000000000000000000}%
\regfield{SYSTEM\_CPU\_INTR\_FROM\_CPU\_3}{1}{0}{{0}}%
\reglabel{Reset}\regnewline%
\begin{regdesc}\begin{reglist}
\label{fielddesc:SYSTEMCPUINTRFROMCPU3}\item [SYSTEM\_CPU\_INTR\_FROM\_CPU\_3] 置 1 生成 CPU 中断 3。该位需在 ISR 过程中由软件清 0。（读/写）
\end{reglist}\end{regdesc}
\end{register}

\begin{register}{H}{SYSTEM\_RSA\_PD\_CTRL\_REG}{\hexprefix{}0068}\label{regdesc:SYSTEMRSAPDCTRLREG}
\regfield{(reserved)}{29}{3}{00000000000000000000000000000}%
\regfield{SYSTEM\_RSA\_MEM\_FORCE\_PD}{1}{2}{{0}}%
\regfield{SYSTEM\_RSA\_MEM\_FORCE\_PU}{1}{1}{{0}}%
\regfield{SYSTEM\_RSA\_MEM\_PD}{1}{0}{{1}}%
\reglabel{Reset}\regnewline%
\begin{regdesc}\begin{reglist}
\label{fielddesc:SYSTEMRSAMEMPD}\item [SYSTEM\_RSA\_MEM\_PD] 置 1 关闭 RSA 内存。该位优先级最低。当 RSA\_DS 外设占用 RSA 加速器时，该位无效。（读/写）
\label{fielddesc:SYSTEMRSAMEMFORCEPU}\item [SYSTEM\_RSA\_MEM\_FORCE\_PU] 置 1 打开 RSA 内存。该位优先级第二高。（读/写）
\label{fielddesc:SYSTEMRSAMEMFORCEPD}\item [SYSTEM\_RSA\_MEM\_FORCE\_PD] 置 1 关闭 RSA 内存。该位优先级最高。（读/写）
\end{reglist}\end{regdesc}
\end{register}


\begin{register}{H}{SYSTEM\_BUSTOEXTMEM\_ENA\_REG}{\hexprefix{}006C}\label{regdesc:SYSTEMBUSTOEXTMEMENAREG}
\regfield{(reserved)}{31}{1}{0000000000000000000000000000000}%
\regfield{SYSTEM\_BUSTOEXTMEM\_ENA}{1}{0}{{1}}%
\reglabel{Reset}\regnewline%
\begin{regdesc}\begin{reglist}
\label{fielddesc:SYSTEMBUSTOEXTMEMENA}\item [SYSTEM\_BUSTOEXTMEM\_ENA] 置 1 使能 EDMA 总线。（读/写）
\end{reglist}\end{regdesc}
\end{register}


\begin{register}{H}{SYSTEM\_CACHE\_CONTROL\_REG}{\hexprefix{}0070}\label{regdesc:SYSTEMCACHECONTROLREG}
\regfield{(reserved)}{29}{3}{00000000000000000000000000000}%
\regfield{SYSTEM\_PRO\_CACHE\_RESET}{1}{2}{{0}}%
\regfield{SYSTEM\_PRO\_DCACHE\_CLK\_ON}{1}{1}{{1}}%
\regfield{SYSTEM\_PRO\_ICACHE\_CLK\_ON}{1}{0}{{1}}%
\reglabel{Reset}\regnewline%
\begin{regdesc}\begin{reglist}
\label{fielddesc:SYSTEMPROICACHECLKON}\item [SYSTEM\_PRO\_ICACHE\_CLK\_ON] 置 1 使能 i-cache 时钟。（读/写）
\label{fielddesc:SYSTEMPRODCACHECLKON}\item [SYSTEM\_PRO\_DCACHE\_CLK\_ON] 置 1 使能 d-cache 时钟。（读/写）
\label{fielddesc:SYSTEMPROCACHERESET}\item [SYSTEM\_PRO\_CACHE\_RESET] 置 1 复位 cache。（读/写）
\end{reglist}\end{regdesc}
\end{register}


\begin{register}{H}{SYSTEM\_EXTERNAL\_DEVICE\_ENCRYPT\_DECRYPT\_CONTROL\_REG}{\hexprefix{}0074}\label{regdesc:SYSTEMEXTERNALDEVICEENCRYPTDECRYPTCONTROLREG}
\regfield{(reserved)}{28}{4}{0000000000000000000000000000}%
\regfield{\tiny SYSTEM\_ENABLE\_DOWNLOAD\_MANUAL\_ENCRYPT}{1}{3}{{0}}%
\regfield{\tiny SYSTEM\_ENABLE\_DOWNLOAD\_G0CB\_DECRYPT}{1}{2}{{0}}%
\regfield{\tiny SYSTEM\_ENABLE\_DOWNLOAD\_DB\_ENCRYPT}{1}{1}{{0}}%
\regfield{\tiny SYSTEM\_ENABLE\_SPI\_MANUAL\_ENCRYPT}{1}{0}{{0}}%
\reglabel{Reset}\regnewline%
\begin{regdesc}\begin{reglist}
\label{fielddesc:SYSTEMENABLESPIMANUALENCRYPT}\item [SYSTEM\_ENABLE\_SPI\_MANUAL\_ENCRYPT] 置 1 在 SPI Boot 模式下使能手动加密 (Manual Encryption)。（读/写）
\label{fielddesc:SYSTEMENABLEDOWNLOADDBENCRYPT}\item [SYSTEM\_ENABLE\_DOWNLOAD\_DB\_ENCRYPT] 置 1 在 Download Boot 模式下使能自动加密 (Auto Encryption)。（读/写）
\label{fielddesc:SYSTEMENABLEDOWNLOADG0CBDECRYPT}\item [SYSTEM\_ENABLE\_DOWNLOAD\_G0CB\_DECRYPT] 置 1 在 Download Boot 模式下使能自动解密 (Auto Decryption)。（读/写）
\label{fielddesc:SYSTEMENABLEDOWNLOADMANUALENCRYPT}\item [SYSTEM\_ENABLE\_DOWNLOAD\_MANUAL\_ENCRYPT] 置 1 在 Download Boot 模式下使能手动加密 (Manual Encryption)。（读/写）
\end{reglist}\end{regdesc}
\end{register}

\begin{register}{H}{SYSTEM\_RTC\_FASTMEM\_CONFIG\_REG}{\hexprefix{}0078}\label{regdesc:SYSTEMRTCFASTMEMCONFIGREG}
\regfield{SYSTEM\_RTC\_MEM\_CRC\_FINISH}{1}{31}{{0}}%
\regfield{SYSTEM\_RTC\_MEM\_CRC\_LEN}{11}{20}{{\hexprefix{}7ff}}%
\regfield{SYSTEM\_RTC\_MEM\_CRC\_ADDR}{11}{9}{{\hexprefix{}0}}%
\regfield{SYSTEM\_RTC\_MEM\_CRC\_START}{1}{8}{{0}}%
\regfield{(reserved)}{8}{0}{00000000}%
\reglabel{Reset}\regnewline%
\begin{regdesc}\begin{reglist}
\label{fielddesc:SYSTEMRTCMEMCRCSTART}\item [SYSTEM\_RTC\_MEM\_CRC\_START] 置 1 启动 RTC 内存的 CRC 校验。（读/写）
\label{fielddesc:SYSTEMRTCMEMCRCADDR}\item [SYSTEM\_RTC\_MEM\_CRC\_ADDR] 设置 CRC 校验的 RTC 存储地址。（读/写）
\label{fielddesc:SYSTEMRTCMEMCRCLEN}\item [SYSTEM\_RTC\_MEM\_CRC\_LEN] 设置用于 CRC 校验的 RTC 存储长度（基于起始地址）。（读/写）
\label{fielddesc:SYSTEMRTCMEMCRCFINISH}\item [SYSTEM\_RTC\_MEM\_CRC\_FINISH] 储存 RTC 存储 CRC 校验状态。高电平表示校验完成，低电平表示校验未完成。（只读）
\end{reglist}\end{regdesc}
\end{register}


\begin{register}{H}{SYSTEM\_RTC\_FASTMEM\_CRC\_REG}{\hexprefix{}007C}\label{regdesc:SYSTEMRTCFASTMEMCRCREG}
\regfield{SYSTEM\_RTC\_MEM\_CRC\_RES}{32}{0}{{0}}%
\reglabel{Reset}\regnewline%
\begin{regdesc}\begin{reglist}
\label{fielddesc:SYSTEMRTCMEMCRCRES}\item [SYSTEM\_RTC\_MEM\_CRC\_RES] 储存 RTC 存储的 CRC 校验结果。（只读）
\end{reglist}\end{regdesc}
\end{register}


\begin{register}{H}{SYSTEM\_SRAM\_CTRL\_2\_REG}{\hexprefix{}0088}\label{regdesc:SYSTEMSRAMCTRL2REG}
\regfield{(reserved)}{10}{22}{0000000000}%
\regfield{SYSTEM\_SRAM\_FORCE\_PU}{22}{0}{{\hexprefix{}3fffff}}%
\reglabel{Reset}\regnewline%
\begin{regdesc}\begin{reglist}
\label{fielddesc:SYSTEMSRAMFORCEPU}\item [SYSTEM\_SRAM\_FORCE\_PU] 控制内部 SRAM 的上电。详见表 \ref{table:sramcontrollingbit}。（读/写）

\end{reglist}\end{regdesc}
\end{register}


\begin{register}{H}{SYSTEM\_SYSCLK\_CONF\_REG}{\hexprefix{}008C}\label{regdesc:SYSTEMSYSCLKCONFREG}
\regfield{(reserved)}{13}{19}{000000000000}%
\regfield{SYSTEM\_CLK\_XTAL\_FREQ}{7}{12}{{0}}%
\regfield{SYSTEM\_SOC\_CLK\_SEL}{2}{10}{{0}}%
\regfield{SYSTEM\_PRE\_DIV\_CNT}{10}{0}{{\hexprefix{}1}}%
\reglabel{Reset}\regnewline%
\begin{regdesc}\begin{reglist}
\label{fielddesc:SYSTEMPREDIVCNT}\item [SYSTEM\_PRE\_DIV\_CNT] 设置预分频器计数器。具体配置，请见章节 \ref{mod:resclk} \textit{\nameref{mod:resclk}}中表 \ref{table:peri clk owner}。（读/写）
\label{fielddesc:SYSTEMSOCCLKSEL}\item [SYSTEM\_SOC\_CLK\_SEL] 选择 SoC 时钟。具体配置，请见章节 \ref{mod:resclk} \textit{\nameref{mod:resclk}}中表 \ref{table:cpu clksrc}。（读/写）
\label{fielddesc:SYSTEMCLKXTALFREQ}\item [SYSTEM\_CLK\_XTAL\_FREQ] 读取晶振频率（单位：MHz）。（只读）
\end{reglist}\end{regdesc}
\end{register}


\begin{register}{H}{SYSTEM\_DATE\_REG}{\hexprefix{}0FFC}\label{regdesc:SYSTEMREGDATEREG}
\regfield{(reserved)}{4}{28}{0000}%
\regfield{SYSTEM\_REG\_DATE}{28}{0}{{\hexprefix{}1908020}}%
\reglabel{Reset}\regnewline%
\begin{regdesc}\begin{reglist}
\label{fielddesc:SYSTEMSYSTEMREGDATE}\item [SYSTEM\_DATE] 版本控制寄存器。（读/写）
\end{reglist}\end{regdesc}
\end{register}
